\documentclass{jsarticle}
\usepackage[dvipdfmx]{graphicx} % Required for inserting images
\usepackage[dvipdfmx]{color}
\usepackage{bm}
\usepackage{amsmath}
\usepackage{here}
\usepackage{url}
\usepackage{siunitx}
\usepackage[top=40truemm,bottom=40truemm,left=30truemm,right=30truemm]{geometry}
\usepackage{fancyhdr}
\usepackage{multirow}
\usepackage{comment}
\usepackage{physics2}

\usepackage[dvipdfmx, hidelinks]{hyperref}

\begin{document}


%//\makeatletter
%//\newcommand*{\themonth}{\two@digits\month}
%//\newcommand*{\theday}{\two@digits\day}
%//\makeatother
%//\renewcommand{\today}{{\the\year}年{\themonth}月{\thetoday}日}
%//\makeatletter

%\newcommand{\todayAD}{\number\year 年\number\month  月\number\day 日}
%↑1桁表示

\newcommand{\todayAD}{
  \number\year 年
  \ifnum\month<10 0\fi\number\month 月
  \ifnum\day<10 0\fi\number\day 日}
%\ifnum\month<10 0\fi は、月が10未満（1桁）のときに0を追加
%\number\month は月の数字を出力する
%同様に、日も \ifnum\day<10 0\fi\number\day で処理
\newcommand{\dy}{$\rm Dy_2Pt_3Al_4$}
\newcommand{\gd}{$\rm Gd_2Pt_3Al_4$}
\newcommand{\ho}{$\rm Ho_2Pt_3Al_4$}
\newcommand{\dys}[1]{$\rm Dy_2Pt_3Al_4 \,\#s{#1}$}
\newcommand{\dyp}[1]{$\rm Dy_2Pt_3Al_4 \,\#p{#1}$}
\renewcommand{\contentsname}{目次}
\renewcommand{\figurename}{図}
\renewcommand{\tablename}{表}
\renewcommand{\refname}{参考文献}

\newcommand{\SIeq}[2]{\SI[parse-numbers=false]{#1}{#2}}
%数式＋単位のとき

\DeclareSIUnit\angstrom{\text{\AA}}
%Å定義(siunitx非推奨)。これで\angstromが問題なく使える。

\begin{flushright}
    \todayAD
\end{flushright}
\begin{center}
    令和7年度\\
    学士論文\\
    \vskip\baselineskip
    新規低対称・多サイト希土類カゴ状化合物\dy の単結晶育成及び物性測定
    \vskip\baselineskip
    22RP019 小林歩夢
    \vskip\baselineskip
    指導教員\\
    道村真司 准教授
    \vskip\baselineskip
    埼玉大学 理学部 物理学科\\
    小坂・道村・佐藤研究室
\end{center}

\newpage
\tableofcontents
\newpage

\section{序論}


\bibliography{myrefs}
\bibliographystyle{junsrt_3}


\begin{comment}
\begin{table}[H]
\begin{center}
\caption{MiniFlex600 測定条件}
\label{t:mf600zyoken}
\begin{tabular}{c|l}
    \hline
    \multicolumn{2}{c}{MiniFlex600 測定条件} \\ \hline\hline
    X線源 & Cu 管球 \SI{40}{kV} \SI{15}{mA} \\ \hline
    ゴニオメーター半径 & \SI{150}{mm} \\ \hline
    スリット条件 & 可変＋固定スリットシステム\\ \hline
    光学素子 & Soller(inc.) : \SI{2.5}{\degree} \\ \hline
    & IHS : \SI{10.0}{mm} \\ \hline
    & DS : \SI{0.625}{\degree} \\ \hline
    & SS : \SI{13.0}{mm}(Open) \\ \hline
    & Soller(rec.) : \SI{2.5}{\degree} \\ \hline
    & RS1 : \SI{8.0}{mm} \\ \hline
    & RS2 : \SI{13.0}{mm}(Open) \\ \hline
    & 単色化 : ${\rm K}\beta$フィルター(×1)[Ni箔] \\ \hline
    試料台 & 標準ガラス板(裏面) \\ \hline
    検出器 & 高速1次元検出器 D/teX Ultra2 \\ \hline
    スキャン条件 & スキャン軸 : $2\theta /\theta$ \\ \hline
    & モード : 連続 \\ \hline
    & 範囲($2\theta$) : \SIeq{3 - 120}{\degree} \\ \hline
    & ステップ : \SI{0.01}{deg/step} \\ \hline
    & スピード : \SI{3.0}{deg/min} \\ \hline


\end{tabular}
\end{center}
\end{table}
\end{comment}

\end{document}